% =====================================================================
%  CHAPTER 8: 社会医学研究方法（对应大纲第八章）
%  参考PPT：第八章 社会医学研究方法
% =====================================================================
\chapter{社会医学研究方法}
\setcounter{chapter}{8}
\chapterintro{
掌握定性研究的概念、特点和常用方法（深入访谈、焦点小组讨论）；掌握问卷的概念、类型、结构和设计原则；掌握信度和效度的概念及类型（难点*）；掌握信度和效度的关系（难点*）；熟悉定量研究的概念和特点（访谈法/自填法的优缺点）；熟悉非概率抽样方法（偶遇/立意/雪球/定额）；熟悉社会医学研究的步骤和课题评价三原则。
}

\section{社会医学研究方法概述}

\subsection{相关研究方法}

\begin{defbox}
社会医学的相关研究方法：调查研究、实验研究（现场实验/社区干预实验）、评价研究（生命质量评价/健康危险因素评价/卫生服务评价/卫生项目评价）、德尔菲法（核心是通过多轮函询征求专家对某一问题的意见）、文献研究（利用已有资料通过综合整理/分析/归纳达到研究目的）。
\end{defbox}

\subsection{社会医学研究的步骤}

\begin{defbox}
第一步选择课题陈述假设→第二步制定研究方案→第三步收集资料→第四步整理和分析资料→第五步解释结果。

\textbf{课题评价三原则：}需要性原则、创造性原则、科学性原则。还需进行可行性论证（从技术和经济角度论证课题是否具备进行研究的条件）。
\end{defbox}

\subsection{非概率抽样方法}

\begin{keybox}
\begin{enumerate}[leftmargin=*]
    \item \imp{偶遇抽样（Accidental/Convenience Sampling）/ 便利抽样}：研究者根据现实情况使用最便利的方式选取样本，简单易行但对总体代表性差
    \item \imp{立意抽样（Purposive Sampling）/ 判断抽样}：调查者根据研究目的分析判断来选择调查对象，与研究目的有关，研究者的判断更为重要
    \item \imp{雪球抽样（Snowball Sampling）}：从总体中少数成员入手，请他们介绍所认识的其他符合条件的人，如同滚雪球一样找到越来越多具有相同性质的群体成员
    \item \imp{定额抽样（Quota Sampling）/ 配额抽样}：依据可能影响研究指标的各种因素对总体分层，确定各层样本比例，再在各层中抽取样本（可看成是分层抽样的延伸但不是随机抽样）
\end{enumerate}
\end{keybox}

\section{定量研究}

\begin{defbox}
\textbf{定量研究（Quantitative Research）/ 量化研究：}通过调查收集人群发生某种事件的数量指标，如患病率、就诊率等，或者探讨各种因素与疾病、健康间的数量依存关系的研究。

\textbf{特点：}优点——重点在研究假设，逻辑推理严谨；标准化和精确化程度高；可用概率统计方法检验；具有较好的客观性和科学性。局限——投入较大；难获得深层次或意料之外信息；一些社会因素及医学问题难以用数据指标表达；一些关系很难用定量结果解释。

\textbf{收集资料方法：}
\begin{itemize}[leftmargin=*]
    \item \imp{访谈法（面对面访谈）}——优点：灵活、适用范围广（含文盲者）、回收率高、可判断真实性、可防第三者影响、可列入复杂问题。缺点：组织工作量大、耗费时间和人力物力、受调查员影响大可能出现访谈偏误、一般没有匿名保证
    \item \imp{自填法}——主要有信访法、现场自填法。信访法优点：节省人力财力时间、有较高匿名保证、地理范围广。信访法缺点：回收率低、被调查者遇到问题无法得到解答、无法控制填写环境、无法判断回答真实性
\end{itemize}
\end{defbox}

\section{问卷设计}

\begin{keybox}
\textbf{问卷（Questionnaire）/ 调查表：}指有详细问题和可供选择的答案或只有问题而无答案的可用来提问收集资料的工具。

\textbf{问卷结构：}封面信（调查目的/内容/意义/请求合作）→指导语（填答说明/概念解释）→问题及答案（特征问题/行为问题/态度问题）→编码（问卷编码/问题编码/答案编码）。

\textbf{问题设计：}
\begin{itemize}[leftmargin=*]
    \item \imp{开放式问题}：只有提问，没有备选答案，被调查者可自由回答。优点：可用于未知答案情况、能收集更多反映特征资料、灵活不受限制、答案过多(>10种)时适用。缺点：可能得到无关信息、要求一定文化水平/语言表达能力、费时费力、应答率低、统计处理困难
    \item \imp{封闭式问题}：在提问同时提供两个及以上备选答案，遵循独立原则和穷尽原则，一般要使至少80\%回答者能找到自己的答案。优点：易回答省时、拒答率低应答率高、低文化程度也能答、敏感问题用等级资料更能获得真实回答、便于统计分析。缺点：有猜答和随便回答可能、答案不易穷尽、了解难以深入、容易笔误
\end{itemize}

\textbf{答案格式种类：}填空式、二项选择式、多项选择式、图表式（线性尺度和表格）、排序式。

\textbf{问题排列原则：}先易后难、封闭式在前开放式在后、问题排列要有逻辑顺序、检验信度的配对问题须分隔开来。

\textbf{问卷设计原则：}目的原则（每个问题都应与研究目的相关）、实用性原则（用词简明清楚准确、避免专业术语、回忆时间不宜过长）、反向性原则（从最终想要结果反推问题）。

\textbf{问卷设计步骤：}明确研究目的→建立问题库（头脑风暴法/借用其它问卷条目）→设计问卷初稿→试用和修改→信度与效度检验。
\end{keybox}

\section{信度与效度（难点*）}

\begin{keybox}
\textbf{信度（Reliability）：}指测量工具的稳定性，表示测量工具（主要指问卷）所获资料的可靠程度或可信程度，即测量的稳定性和一致性问题。

\textbf{信度测量方式：}
\begin{enumerate}[leftmargin=*]
    \item \imp{复测信度（Test-Retest Reliability）}：采用同一问卷在同一人群中测量两次，评价两次测量的相关性
    \item \imp{复本信度（Alternate Form Reliability）}：设计另一种与研究问卷高度类似的问卷，同时测量研究对象
    \item \imp{折半信度（Split-Half Reliability）}：将一个问卷分拆为两半，分别作为各自的复本
\end{enumerate}

\textbf{效度（Validity）：}调查结果与试图要达到的目标之间的接近程度，即调查的结果是否是需要的结果。评价的是偏倚问题，测量的真实性问题。

\textbf{效度的种类：}
\begin{enumerate}[leftmargin=*]
    \item \imp{表面效度（Face Validity）}：从表面上看问卷条目是否都与研究者想了解的问题有关
    \item \imp{内容效度（Content Validity）}：评价问卷所涉及内容能在多大程度上覆盖研究目的达到的各个方面和领域（由专家进行主观评价）
    \item \imp{结构效度（Construct Validity）}：一个概念或命题按一定理论分解后形成相应结构，测量的实际结构与此一致的程度
    \item \imp{准则效度（Criterion Validity）}：问卷结构与某个"金标准"的一致性，即问卷调查结果与标准测量结果间的接近程度
\end{enumerate}

\textbf{信度和效度的关系（难点*）：}
\begin{enumerate}[leftmargin=*]
    \item 不可信的测量一定是无效的（信度不高，效度也不会高）
    \item 可信的测量既可能有效，也可能无效（信度高，效度不一定高）
    \item 无效的测量既可能是可信的，也可能是不可信的（效度不高，信度可能高，也可能不高）
    \item \imp{有效的测量一定是可信的测量（效度高，信度一定也高）}
\end{enumerate}
\end{keybox}

\section{定性研究}

\begin{keybox}
\textbf{定性研究（Qualitative Research）/ 质性研究：}是一种在自然的情境下，从整体的角度深入探讨和阐述被研究事物的特点及其发生和发展规律，以揭示事物的内在本质的一类研究方法。

\textbf{特点：}注重事物的过程；小样本人群研究；需要与研究对象保持较长时间的接触；研究结果很少用概率统计分析。

\textbf{应用：}(1)用于帮助问卷设计；(2)验证因果关系，探讨发生机制；(3)分析定量研究出现矛盾结果的原因；(4)了解危险因素的变化情况；(5)作为快速评价技术，为其它研究提供信息。

\textbf{常用定性研究方法：}
\begin{enumerate}[leftmargin=*]
    \item \imp{深入访谈（In-Depth Interview）}：调查者与回答者进行一对一深入交谈来收集资料
    \item \imp{焦点小组讨论（Focus Group Discussion）}：由一位主持人带领，围绕某一主题进行自由、深入讨论（一般6$\sim$10人，1$\sim$2小时），了解调查对象对某一问题的看法
    \item \imp{名义小组讨论（Nominal Group Discussion）}
    \item \imp{观察法（Observation）}
\end{enumerate}
\end{keybox}

\section{复习题}

\begin{enumerate}[leftmargin=*, itemsep=8pt]
    \item \textbf{【名词解释】}定量研究；定性研究；问卷；信度；效度；德尔菲法
    \item \textbf{【名词解释】}偶遇抽样；立意抽样；雪球抽样；定额抽样
    \item \textbf{【名词解释】}复测信度；表面效度；内容效度；准则效度；结构效度
    \item \textbf{【简答题】}简述课题评价三原则和非概率抽样的四种方法。
    \item \textbf{【简答题】}简述面对面访谈法和信访法的优缺点。
    \item \textbf{【简答题】}简述开放式问题和封闭式问题的定义及其优缺点。
    \item \textbf{【简答题】}简述问卷的结构、设计原则和步骤。
    \item \textbf{【简答题】}简述信度和效度的概念、类型及其相互关系（难点*）。
    \item \textbf{【简答题】}简述定性研究的概念、特点、应用和常用方法。
\end{enumerate}

\subsection*{参考答案要点}

\begin{enumerate}[leftmargin=*, itemsep=5pt]
    \item \textbf{定量研究}：通过调查收集人群发生某种事件的数量指标，探讨各种因素与疾病/健康间数量依存关系的研究。\textbf{定性研究}：在自然情境下从整体角度深入探讨和阐述被研究事物的特点及规律、揭示事物内在本质的研究方法。\textbf{问卷}：有详细问题和可供选择的答案或只有问题而无答案的可用来收集资料的工具。\textbf{信度}：测量工具的稳定性，表示所获资料的可靠程度。\textbf{效度}：调查结果与试图达到目标之间的接近程度，评价偏倚和真实性。\textbf{德尔菲法}：通过多轮函询征求专家意见的方法。
    \item 偶遇抽样（最便利方式选取，代表性差）、立意抽样（根据研究目的分析判断选择，研究者的判断更为重要）、雪球抽样（从少数成员入手，滚雪球式扩大）、定额抽样（先分层确定比例再在各层抽取，非随机）。
    \item 复测信度：同一问卷同一人群测两次评价相关性。表面效度：问卷条目从表面看是否与研究问题有关。内容效度：问卷内容覆盖研究目的各方面和领域的程度（专家主观评价）。准则效度：问卷结构与"金标准"的一致性。结构效度：概念按理论分解后测量的实际结构与之的一致程度。
    \item 三原则：需要性、创造性、科学性+可行性论证。四种非概率抽样见第2题。
    \item 面对面访谈优点：灵活、适用范围广、回收率高、可判断真实性、防第三者影响。缺点：组织工作量大、费时费力、受调查员影响大、无匿名保证。信访法优点：省人力财力时间、匿名保证高、地理范围广。缺点：回收率低、无法解答疑问、无法控制环境、无法判断真实性。
    \item 开放式问题：只有提问无备选答案，自由回答。优点：未知答案情况适用、能收集更多资料、灵活。缺点：可能得无关信息、需文化水平、费时、应答率低、统计困难。封闭式问题：提供备选答案。优点：易答省时、应答率高、低文化可答、敏感问题更真实、便于统计。缺点：可能猜答、答案不易穷尽、难以深入了解、易笔误。
    \item 结构：封面信→指导语→问题及答案→编码。设计原则：目的原则、实用性原则、反向性原则。步骤：明确目的→建立问题库→设计初稿→试用修改→信度效度检验。
    \item 信度三类型：复测信度、复本信度、折半信度。效度四类型：表面效度、内容效度、结构效度、准则效度。信效度关系：(1)不可信必无效；(2)可信可能有效也可能无效；(3)无效可能可信也可能不可信；(4)有效必可信（效度高信度一定高）。
    \item 特点：注重过程、小样本人群、长时间接触、很少用概率统计。应用：帮助问卷设计、验证因果关系、分析矛盾结果、了解危险因素变化、快速评价。常用方法：深入访谈、焦点小组讨论（6$\sim$10人，1$\sim$2小时）、名义小组讨论、观察法。
\end{enumerate}

\newpage
